\documentclass[a4paper, 12pt,french]{report}
%====================== PACKAGES ======================
\usepackage[utf8]{inputenc}
\usepackage[T1]{fontenc}
\usepackage[french]{babel}
%pour gerer les positionnement d'images
\usepackage{float}
\usepackage{amsmath}
\usepackage{graphicx}
\usepackage[colorinlistoftodos]{todonotes}
\usepackage{url}
%pour les informations sur un document compilé en PDF et les liens
%externes / internes 
\usepackage{hyperref}
%pour la mise en page des tableaux
\usepackage{array}
\usepackage{tabularx}
%espacement entre les lignes
\usepackage{setspace}
%modifier la mise en page de l'abstract
\usepackage{abstract}
%police et mise en page (marges) du document
\usepackage[T1]{fontenc}
\usepackage[top=2cm, bottom=2cm, left=2cm, right=2cm]{geometry}
%Pour les galerie d'images
\usepackage{subfig}
%Pour les accents
\usepackage{amssymb}
\usepackage{amsthm}
\usepackage{stmaryrd}
\usepackage{dsfont}
\usepackage{mathtools}
\usepackage{listings}
\usepackage{hyperref}
\usepackage{algorithm}
\usepackage{algpseudocode}
\usepackage{relsize}
\DeclareMathOperator{\Tr}{Tr}

\title{HMM}
\author{Nonolewawa }
\date{January 2025}

\begin{document}

\maketitle

\section{Sliced Score matching}
\noindent
Soit $\nu$ un vecteur aléatoire à valeur dans $\mathbb{R}^n$ tel que $\mathbb{E}_\nu[\nu_i\nu_j]=\delta_{ij}$. Nous avons le résultat suivant : 

\begin{align*}
\mathcal{L}(\theta) & = \int_0^T\lambda(t)\mathbb{E}_{X_t}[||s_\theta(X_t,t)-\nabla_{X_t}\log p_t(X_t)||^2]dt \\
& = \int_0^T\lambda(t)\mathbb{E}_{X_t}[||s_\theta(X_t,t)||^2 + 
2\mathbb{E}_\nu[\frac{d}{dh}\nu^Ts_\theta(X_t+h\nu,t)|_{h=0}]
]dt + C
\end{align*}
Où $C$ est une constante par rapport à $\theta$. Nous montrerons ce résultat en trois étapes  :

%Mettre dans le préambule \DeclareMathOperator{\Tr}{Tr}
\begin{enumerate}
    \item $\forall A \in \mathcal{M}_n(\mathbb{R}), \mathbb{E}_\nu [\nu^TA\nu]=\Tr(A)$
    \item $\mathcal{L}(\theta)=\int_0^T\lambda(t)\mathbb{E}_{X_t}[||s_\theta(X_t,t)||^2 -2 
    \langle s_\theta(X_t,t), \nabla_{X_t}\log p_t(X_t) \rangle
]dt + C$
    \item $-\mathbb{E}_{X_t}[\langle s_\theta(X_t,t), \nabla_{X_t}\log p_t(X_t) \rangle] =
    \mathbb{E}_{X_t}\mathbb{E}_\nu[\frac{d}{dh}\nu^Ts_\theta(X_t+h\nu,t)|_{h=0}]$
\end{enumerate}

\noindent
1. On note $A=(a_{ij})$. \\
$$\mathbb{E}_\nu[\nu^TA\nu]
=\mathbb{E}_\nu[\sum_{i,j}a_{ij}\nu_i\nu_j]
=\sum_{i,j}a_{ij}\mathbb{E}_\nu[\nu_i\nu_j]
=\sum_ia_{ii}
=\Tr(A)$$
2. 
$$||s_\theta(X_t,t)-\nabla_{X_t}\log p_t(X_t)||^2=||s_\theta(X_t,t)||^2
-2\langle s_\theta(X_t,t), \nabla_{X_t}\log p_t(X_t) \rangle
+||\nabla_{X_t}\log p_t(X_t)||^2$$
Comme $\int_0^T\lambda(t)\mathbb{E}_{X_t}[||\nabla_{X_t}\log p_t(X_t)||^2]dt$ est constant par rapport à $\theta$, on peut noter ce terme par $C$, ce qui donne le résultat. \\
\noindent
3.
\begin{align*}
-\mathbb{E}_{X_t}[\langle s_\theta(X_t,t), \nabla_{X_t}\log p_t(X_t) \rangle] &
=-\int_{\mathbb{R}^n}\langle s_\theta(x,t), \nabla_{x}\log p_t(x) \rangle p_t(x)dx \\
& = -\int_{\mathbb{R}^n}\langle s_\theta(x,t), \nabla_{x} p_t(x) \rangle dx
\end{align*}
Car $\nabla_{x}\log p_t(x)=\nabla_{x}p_t(x) / p_t(x)$. \\
On note $s_\theta(x,t)=(s_\theta(x,t)_1,s_\theta(x,t)_2,...,s_\theta(x,t)_n)$.
Pour $i$ compris entre $1$ et $n$, on a par intégration par partie : 
$$\int_{x_i \in \mathbb{R}}s_\theta(x,t)_i\frac{\partial pt(x)}{\partial x_i}dx_i=
[s_\theta(x,t)_ipt(x)]_{x_i=-\infty}^{x_i=+\infty}
-\int_{x_i \in \mathbb{R}}\frac{\partial s_\theta(x,t)_i}{\partial x_i}pt(x)dx_i$$
Ce qui, avec le théorème de Fubini, donne : 
\begin{align*}
\int_{\mathbb{R}^n}s_\theta(x,t)_i\frac{\partial pt(x)}{\partial x_i}dx & =
\int_{x_1 \in \mathbb{R}}\int_{x_2 \in \mathbb{R}}...\int_{x_n \in \mathbb{R}}s_\theta(x,t)_i\frac{\partial pt(x)}{\partial x_i}dx_n...dx_2dx_1 \\
& = \int_{x_1 \in \mathbb{R}}...\int_{x_n \in \mathbb{R}}\int_{x_i \in \mathbb{R}}s_\theta(x,t)_i\frac{\partial pt(x)}{\partial x_i}dx_idx_n...dx_1 \\
& = -\int_{x_1 \in \mathbb{R}}...\int_{x_n \in \mathbb{R}}\int_{x_i \in \mathbb{R}}\frac{\partial s_\theta(x,t)_i}{\partial x_i}pt(x)dx_idx_n...dx_1 \\
& = -\int_{\mathbb{R}^n}\frac{\partial s_\theta(x,t)_i}{\partial x_i}pt(x)dx
\end{align*}
Ainsi, par linéarité de l'intégrale et en appliquant 1., on obtient : 
\begin{align*}
-\mathbb{E}_{X_t}[\langle s_\theta(X_t,t), \nabla_{X_t}\log p_t(X_t) \rangle] &=
\int_{\mathbb{R}^n}\sum_i\frac{\partial s_\theta(x,t)_i}{\partial x_i}pt(x)dx \\
&= \mathbb{E}_{X_t}[\Tr(J_{X_t}s_\theta(X_t,t))] \\
&= \mathbb{E}_{X_t}\mathbb{E}_\nu[\nu^TJ_{X_t}s_\theta(X_t,t)\nu]
\end{align*}
Où $J_{X_t}s_\theta(X_t,t)$ est la matrice Jacobienne de $x \mapsto s_\theta(x,t)$ au point $X_t$.
Pour la dernière étape, on pose : 
$$ 
\eta : 
\left \{
\begin{array}{rcl}
\mathbb{R} & \longrightarrow & \mathbb{R}^n  \\
h & \longmapsto & s_\theta(X_t+h\nu,t) 
\end{array}
\right.
$$
On a : 
$$ \frac{d}{dh}s_\theta(X_t+h\nu,t)=\eta'(h)
=J_{X_t+h\nu}s_\theta(X_t+h\nu,t)\nu
$$
Et ainsi : 
$$
\nu^TJ_{X_t}s_\theta(X_t,t)\nu=
\nu^TJ_{X_t+h\nu}s_\theta(X_t+h\nu,t)\nu|_{h=0}=
\frac{d}{dh}\nu^Ts_\theta(X_t+h\nu,t)|_{h=0}
$$
Ce qui donne le résultat de l'étape 3. \\
Pour déduire la formule du Sliced Score Matching, il suffit de d'injecter l'égalité de l'étape 3. dans l'étape 2. ce qui donne la nouvelle formule de $\mathcal{L}(\theta)$ montrée au début de cette partie.



\end{document}